\section{The residue action and its four eta products}
This section identifies the modular forms arising from the cyclic action on
four copies of the root lattice $A_5$.
Let $N$ be the Niemeier lattice with root system $A_5^4D_4$.
Let $c$ cyclically permute the six coordinates of each $A_5$ factor and
act trivially on $D_4$.  Root reflections act trivially on the
root-lattice discriminant group, so $c$ preserves the glue defining $N$.
A Frame shape $\prod_d d^{a_d}$ means the identity
$\det(t-c)=\prod_d(t^d-1)^{a_d}$.

\begin{theorem}\label{thm:residue-eta}
The squares in $(\mathbb Z/840\mathbb Z)^\times$ form the cyclic group
$S=\langle289\rangle$ of order six. Under $289\mapsto c$, an element
of order $d$ has Frame shape $d^{24/d}$. Its eta product is
$F_d(\tau)=\eta(d\tau)^{24/d}$, and the four products are normalized
newforms, as follows:
\[
\begin{array}{c|c|c|c}
 d&F_d&\text{weight}&\text{level}\\\hline
 1&\Delta&12&1\\
 2&\eta(2\tau)^{12}&6&4\\
 3&\eta(3\tau)^8&4&9\\
 6&\eta(6\tau)^4&2&36
\end{array}
\]
Exactly the last two have complex multiplication, in both cases by
$\mathbb Q(\sqrt{-3})$. This occurs exactly when the cubic character
modulo seven is nontrivial on the element of $S$.
\end{theorem}
\begin{proof}
The Chinese remainder theorem gives the group of squares modulo
$840=8\cdot3\cdot5\cdot7$ as the product of groups of orders
$1,1,2,3$. The successive powers of $289$ are
$1,289,361,169,121,529$, and the sixth power is one. The action of a
six-cycle on the sum-zero subspace of $\mathbb Q^6$ has characteristic
polynomial $(t^6-1)/(t-1)$. Four copies, together with the four fixed
coordinates from $D_4$, therefore have polynomial $(t^6-1)^4$.
Equivalently, the rational representation is four copies of the regular
representation of the cyclic group of order six. Restriction to an
element of order $d$ gives $24/d$ regular representations of order $d$.

For $N=d^2$ the eta-product congruences are
$\sum_{\delta\mid N}\delta r_\delta=24$ and
$\sum_{\delta\mid N}(N/\delta)r_\delta=24$, with $r_d=24/d$.
The eta transformation formula gives a modular form with trivial
character; all cusp orders are positive because its only eta exponent
is positive. Its first term is $q$. The spaces
$S_{12}(\mathrm{SL}_2(\mathbb Z))$, $S_6(\Gamma_0(4))$,
$S_4(\Gamma_0(9))$, and $S_2(\Gamma_0(36))$ each have dimension one.
Their lower-level old spaces are zero, so these normalized cusp forms
are newforms and simultaneous Hecke eigenforms.

The complex multiplication assertions for $d=3,6$ also follow from the
explicit characters in Theorem~\ref{thm:primitive-cube}.
A complex multiplication field for a newform has discriminant dividing
its level. No imaginary quadratic discriminant divides one. At level
four the only possibility is $\mathbb Q(i)$, whereas the coefficient
of $q^3$ in $\eta(2\tau)^{12}$ is $-12$; it would vanish at the inert
prime three for a form with this complex multiplication. Thus the first
two forms have no complex multiplication. Finally, if $r=289^j$, its
cubic character modulo seven is trivial exactly when $3\mid j$.
Those are precisely the elements of orders one and two.
\end{proof}

\begin{proposition}[The two residue symbols]\label{prop:residue-symbol}
Put $\omega=e^{2\pi i/3}$, choose the quartic character modulo five
by $(2/5)_4=i$, and put
$\pi=3+\omega^2=2-\omega$, a prime of norm seven in $\mathbb Z[\omega]$.
Then the character of $S$ which sends $289$ to $e^{\pi i/3}$ is
\[
 \kappa(r)=\left(\frac r5\right)_4
                \left(\frac r\pi\right)_3 .
\]
In particular its cubic factor is the restriction to rational integers
of the cubic residue character at $\pi$.
\end{proposition}
\begin{proof}
In $\mathbb Z[\omega]/(\pi)=\mathbb F_7$ one has $\omega=2$,
so $(r/\pi)_3$ is the unique cube root of unity congruent to
$r^{(7-1)/3}=r^2$ modulo $\pi$. At $r=289$ the quartic factor is
$-1$, while $r=2$ modulo seven gives the cubic factor $\omega^2$.
Their product is $-\omega^2=e^{\pi i/3}$. Both factors are
characters; equality at the generator proves the identity on $S$.
\end{proof}

\section{The primitive cubic character and all local factors}
We distinguish a power of a primitive Hecke character from the
imprimitive ideal series obtained by retaining its old excluded primes.
Put $\omega=e^{2\pi i/3}$, $\mathcal O=\mathbb Z[\omega]$, and
$\mathfrak p=(\sqrt{-3})$. The ring $\mathcal O$ is a principal ideal
domain. Every ideal prime to $2\mathfrak p$ has a unique generator
$\alpha\equiv1\pmod{2\mathfrak p}$; every ideal prime to $3$ has a
unique generator $\beta\equiv1\pmod3$. Define their respective
Hecke characters by $\psi_{36}(\mathfrak a)=\alpha$ and
$\psi_{27}(\mathfrak a)=\beta$.

\begin{theorem}\label{thm:primitive-cube}
The characters $\psi_{36}$ and $\psi_{27}$ have infinity type one
and conductors $2\mathfrak p$ and $(3)$. Their primitive cubes coincide:
\[
 \psi_{36}^3=\psi_{27}^3=\Psi,
 \qquad \mathfrak f_\Psi=\mathfrak p.
\]
For an ideal prime to $3$, choose any generator
$\gamma\equiv1\pmod{\mathfrak p}$. Then
$\Psi(\mathfrak a)=\gamma^3$ is well-defined, and
\[
 f=\theta_{\psi_{36}}=\eta(6\tau)^4,
 \quad g=\theta_\Psi=\eta(3\tau)^8,
 \quad \theta_{\psi_{27}}=\eta(3\tau)^2\eta(9\tau)^2.
\]
For the usual arithmetic normalization of the finite local factors,
\begin{equation}\label{eq:full-cube}
 L(\operatorname{Sym}^3f,s)=L(g,s)L(f,s-1)
\end{equation}
with no omitted primes. At $p=2$ the factor on either side is
$(1+8\,2^{-2s})^{-1}$; at $p=3$ it is one. At every prime $p\ge5$,
\[
 a_p(g)=a_p(f)^3-3p\,a_p(f).
\]
\end{theorem}
\begin{proof}
Reduction identifies $\mathcal O^\times$ with
$(\mathcal O/(2\mathfrak p))^\times$ and with
$(\mathcal O/(3))^\times$: all three groups have order six, and the
six units have distinct reductions. This proves the two uniqueness
assertions and multiplicativity. Modulo $\mathfrak p$, the units with
value one are exactly $1,\omega,\omega^2$. Their cubes are one,
so the displayed definition of $\Psi$ is independent of the
choice. The primary generators for both type-one characters are
primary modulo $\mathfrak p$, proving the cube identities away from
six. They therefore identify the primitive characters.

The finite part of a type-one character cannot be trivial on any unit
other than one. Removing either prime factor from $2\mathfrak p$
would make $-1$ or $\omega$ congruent to one; replacing $(3)$ by
$\mathfrak p$ would make $\omega$ congruent to one. Their conductors
are therefore exactly as stated. For type three, modulus one is
impossible because $(-1)^3=-1$, while modulus $\mathfrak p$ works.
The theta identities are the explicit formulas of
Huber, Liu, McLaughlin, Ye, Yuan and Zhang
\cite[Lemma~2.1, equations~(2.2), (2.4), and~(2.10)]{HuberCM}.
In the last formula
each ideal has three primary generators modulo $\mathfrak p$, which
accounts for its factor $1/3$.

Let $\rho=\operatorname{Ind}_{G_K}^{G_{\mathbb Q}}\psi_{36,\ell}$
be the compatible two-dimensional representation of $f$, and let
$\varepsilon$ be the cyclotomic character. Its determinant is
$\varepsilon$. On restriction to $G_K$, write a basis as $e_1,e_2$
with characters $\psi,\psi^c$. In the symmetric cube, the spans
$\langle e_1^3,e_2^3\rangle$ and
$\langle e_1^2e_2,e_1e_2^2\rangle$ are invariant. The first is
$\operatorname{Ind}\psi^3$. The second restricts to
$\psi\varepsilon\oplus\psi^c\varepsilon$ and is isomorphic to
$\rho\otimes\varepsilon$; an apparent quadratic sign in a choice of
induction basis is removed by the standard isomorphism
$\operatorname{Ind}\psi\simeq
 (\operatorname{Ind}\psi)\otimes\chi_K$.
Thus
\[
 \operatorname{Sym}^3\rho\simeq
 \operatorname{Ind}\Psi\oplus(\rho\otimes\varepsilon).
\]
The local Euler factor is the inverse characteristic polynomial of
Frobenius on inertia invariants. Taking inertia invariants respects
direct sums, so this representation identity proves
\eqref{eq:full-cube} at every finite prime, including ramified primes.

At two, $f$ has trivial local factor, while two is inert in $K$ and
$\Psi((2))=(-2)^3=-8$; its norm is four. This gives the claimed
factor. At three, both characters have ramified local components
and the induced representations have no inertia invariants, so both
factors are one. At a split unramified prime, the two Frobenius
eigenvalues $\alpha,\beta$ satisfy $\alpha\beta=p$, and
$\alpha^3+\beta^3=(\alpha+\beta)^3-3p(\alpha+\beta)$.
At an inert prime both traces vanish. This proves the coefficient
identity in all cases.
\end{proof}

\section{Theta coordinates of ternary perpendicular sublattices}
The same two small shell counts determine the theta series for every
isotropic plane modulo three.
Write $E=E_{12}=1+C\sum_{n\ge1}\sigma_{11}(n)q^n$, where
$C=65520/691$, and put
\[
 a=\frac{7381}{66430},\qquad b=\frac{59049}{66430}.
\]
Let $F\in S_{12}(\Gamma_0(3))$ be the normalized newform whose
first coefficients are $q+78q^2-243q^3+4036q^4+\cdots$.

\begin{lemma}\label{lem:level-three-basis}
Put $H=\eta(\tau)^{12}\eta(3\tau)^{12}$. Then
\[
 F=\Delta(\tau)+729\Delta(3\tau)+102H,
 \qquad F|_{12}W_3=F.
\]
The forms $E(\tau),E(3\tau),\Delta(\tau),\Delta(3\tau),F(\tau)$
are a basis of $M_{12}(\Gamma_0(3))$.
\end{lemma}
\begin{proof}
The eta transformation law shows that $H$ is a cusp form of weight
twelve and trivial character at level three. Direct multiplication gives
$H=q^2-12q^3+54q^4+\cdots$. The index of $\Gamma_0(3)$ is four,
so the valence bound gives an injective map from $M_{12}(\Gamma_0(3))$
to its coefficients of $q^0,\ldots,q^4$. The five displayed forms with
$H$ in place of $F$ have independent images. Thus they are a basis;
the cusp space has dimension three and its old subspace is spanned by
$\Delta(\tau),\Delta(3\tau)$.

With the normalization
$T_2\sum c_nq^n=\sum(c_{2n}+2^{11}c_{n/2})q^n$,
coefficient comparison through $q^4$ gives
\[
 T_2H=\Delta+729\Delta(3\tau)+78H.
\]
Here $c_{n/2}=0$ if $n$ is odd; the finite comparison uses the product
for $H$ through $q^8$. Both oldforms have $T_2$ eigenvalue $-24$.
The displayed linear combination $F$ has eigenvalue $78$, so it
spans the one-dimensional new subspace and is its normalized newform.
Finally $\Delta|W_3=729\Delta(3\tau)$ and $H|W_3=H$, proving
the Fricke sign. These comparisons are reproduced with exact rational
coefficients in the accompanying coefficient program.
\end{proof}

\begin{theorem}\label{thm:universal-theta-coordinates}
Let $\Lambda$ be the Leech lattice and let $P$ be a totally isotropic
two-dimensional subspace of $\Lambda/3\Lambda$, with quadratic form
$x\mapsto(x,x)/2\pmod3$. Set
\[
 J_P=\{x\in\Lambda:(x,P)\equiv0\pmod3\},
\]\[ 
 k=\#\{x\in J_P:(x,x)=4\},\qquad
 t=\#\{u\in3J_P^\vee:(u,u)=6\}.
\]
Then $J_P$ has determinant $81$, level three, and theta series
\begin{equation}\label{eq:universal-theta}
 \Theta_{J_P}=aE(\tau)+bE(3\tau)
            +x\Delta(\tau)+y\Delta(3\tau)+zF(\tau),
\end{equation}
where
\[
 z=\frac{k-2073aC}{102},\qquad
 x=-aC-z,\qquad
 y=59049t-bC-729z.
\]
Moreover $t\in\{0,2,4,6,8\}$: half of $t$ is the number of
lines of $P$ represented by norm-six vectors of $\Lambda$.
\end{theorem}
\begin{proof}
Unimodularity makes the pairing modulo three nondegenerate, so the
two independent conditions defining $J_P$ give index nine.
Consequently its determinant is $9^2$. If $\widetilde P$ denotes
any lifts of a basis of $P$, then
$J_P^\vee=\Lambda+\tfrac13\mathbb Z\widetilde P$.
Total isotropy implies $3J_P^\vee\subset J_P$ and
$3(u,u)/2\in\mathbb Z$ for every $u\in J_P^\vee$.
The associated discriminant form is nondegenerate and nontrivial,
so its level is exactly three. Poisson summation and the evenness of
$J_P$ give modularity of weight twelve and trivial character on
$\Gamma_0(3)$, with Fricke transformation
\[
 \Theta_{J_P}|_{12}W_3
       =81\,\Theta_{J_P^\vee}(3\tau).
\]
Here the determinant-normalized slash operator is used, so that
\[E|W_3=729E(3\tau),\qquad E(3\tau)|W_3=E/729.\]

The space $M_{12}(\Gamma_0(3))$ has basis
\[E(\tau),E(3\tau),\Delta(\tau),\Delta(3\tau),F(\tau).\]
The newform satisfies $F|W_3=F$. Comparing constant terms at the two
cusps gives $a+b=1$ and $729a+b/729=81$, hence the displayed
constants. Since $J_P$ contains no roots, its $q$ coefficient is
zero, giving $x+z=-aC$. Its $q^2$ coefficient gives
$k=2049aC-24x+78z=2073aC+102z$.
The coefficient of $q$ after applying $W_3$ is on the one hand
$81t$, and on the other hand $(bC+y)/729+z$. This yields $y$.

Finally $3J_P^\vee$ is the preimage of $P$ in $\Lambda$.
There is at most one norm-six vector in each nonzero residue class
modulo $3\Lambda$: if $u,v$ were distinct, their squared distance
would be at least $36$, the minimum of $3\Lambda$, but at most
$(\sqrt6+\sqrt6)^2=24$. The eight nonzero vectors of $P$ form
four pairs of opposite vectors. Opposite classes have the same
norm-six count, proving the last assertion.
\end{proof}

\section{Exact theta coordinates for the 129 cases}
The theta-coordinate formula reduces each lattice calculation to two
small shell enumerations. We construct the lattices from their defining
glue and plane representatives, then enumerate the two shells exactly.

\begin{theorem}[Finite computation on the representatives]
\label{thm:129-theta}
For each of the 129 plane representatives, the reconstructed lattice
is even, positive definite, rootless, of determinant $81$ and level
three. The pairs $(k,t)$ in
Theorem~\ref{thm:universal-theta-coordinates} give precisely the
18 theta series in Table~\ref{tab:18theta}. The last column counts
representatives, not orbit sizes. Complete rational values of all five
coordinates, separately for every representative, are tabulated by
the workbench program \emph{Ternary theta coordinates}.
\end{theorem}

\begin{table}[htbp]
\centering\small
\caption{The 18 theta series among the 129 plane representatives.
The two Eisenstein coordinates are always $a,b$; $x,y$ are recovered
from $k,t,z$ by Theorem~\ref{thm:universal-theta-coordinates}.}
\label{tab:18theta}
\begin{tabular}{rrrr}
\hline
$k$&$t$&$z$&Number of representatives\\\hline
21600&0&$-2916/1241$&28\\
21762&0&$-945/1241$&15\\
21924&0&$1026/1241$&23\\
22086&0&$2997/1241$&25\\
22248&0&$4968/1241$&11\\
22410&0&$6939/1241$&4\\
22572&0&$8910/1241$&7\\
22896&0&$756/73$&2\\
23058&0&$14823/1241$&1\\
23544&0&$20736/1241$&2\\
25488&2&$44388/1241$&1\\
25650&2&$2727/73$&1\\
25812&2&$48330/1241$&2\\
25974&2&$50301/1241$&2\\
26298&2&$54243/1241$&2\\
29700&4&$95634/1241$&1\\
30348&4&$103518/1241$&1\\
33264&6&$138996/1241$&1\\\hline
\end{tabular}
\end{table}

\begin{proof}[Construction and exhaustive finite verification]
The standard integral-coordinate Leech model was reconstructed from
the twelve binary Golay generators, from the code construction. Every one of the 654 stabilizer generators was
then checked to preserve its Gram matrix and to have determinant
$\pm1$.

For 123 representatives the intrinsic roots, in their marked
unimodular lattice $N$, generate a rank-two subgroup of
$(\tfrac13N)/N$. Their integral annihilator
\[
 \{x\in N:(x,r)\in\mathbb Z\text{ for every root }r\}
\]
has index nine and is the required perpendicular lattice. Both marked
models were independently built: the $A_2^{12}$ model by the explicit
neighbour construction below, and the $A_1^{24}$ model by binary Golay
glue on 24 orthogonal roots of norm two. The squared norm of each
root was checked in the reconstructed metric.

The six remaining representatives were treated separately. If a matrix
preserves a plane over $\mathbb F_3$, its restriction to that plane
satisfies $(X^8-1)^2$. Intersecting the kernels of this polynomial
for the generators and deterministic words in them gives
spaces of dimension two or three in these six cases. Enumerating all
planes in the latter spaces and testing isotropy and invariance
recovers a unique plane for representatives 2, 6, 127, 128 and 129. For representative
4 there are two candidates; both give $(k,t)=(23058,0)$. Independently,
the defining intersection construction below gives the same pair.
Thus this ambiguity does not affect the theta table. In particular,
the reconstruction does not silently choose one of the two planes.

Each resulting Gram matrix $G$ was checked to have determinant 81,
even diagonal, and positive principal minors. Also $3G^{-1}$ was
checked to be integral with even diagonal. This proves the asserted
level; it cannot be one for a nonunimodular even lattice. The Gram
matrix of $3J^\vee$ is $9G^{-1}$.

For enumeration, an exact unimodular basis change is verified first.
Writing $D_i$ for the principal minors, $D_0=1$, fraction-free
elimination gives integers $m_{ij}$ with
\[
 x^tGx=\sum_{i=0}^{23}
 \frac{(D_{i+1}x_i+\sum_{j>i}m_{ij}x_j)^2}{D_iD_{i+1}}.
\]
The complete matrix identity is checked, not just its diagonal.
Descending in $i$, the residual recursion is
\[
 R_i=\frac{D_iR_{i+1}-(D_{i+1}x_i+
                 \sum_{j>i}m_{ij}x_j)^2}{D_{i+1}}.
\]
An integer square root gives the closed interval of possible $x_i$;
every integer in that interval is visited, and each division is
checked to be exact. Induction on $i$ proves that the enumeration
includes all boundary vectors and omits no vector of the stated norm.
Norm-four counts in $J$ use a new C++ implementation; norm-six counts
in $3J^\vee$ use a separately written Python-integer recursion.
The certificates record the original and reduced Gram matrices,
basis changes, principal minors, elimination coefficients, shell
counts, and rational theta coordinates. They reproduce all entries
of the table. Theorem~\ref{thm:universal-theta-coordinates} then
identifies the entire theta series, rather than merely a truncated
expansion.
\end{proof}

\paragraph{Scope of the classification check.}
The computation above verifies the lattices and theta coordinates of
the representatives. It is not a new proof that the representatives
form all 129 orbits under $\mathrm{Co}_0$, nor a new proof of their
full stabilizer orders. Using the orders, the three exact
incidence checks are
\[
\begin{array}{rcl}
 \sum_P |\mathrm{Co}_0|/|\mathrm{Co}_{0,P}|
   &=&61549107892647466000,\\
 \sum_P (|\mathrm{Co}_0|/|\mathrm{Co}_{0,P}|)k(P)
   &=&196560(3^{22}-1)(3^{20}-1)/16,\\
 \sum_P (|\mathrm{Co}_0|/|\mathrm{Co}_{0,P}|)t(P)
   &=&16773120(3^{10}+1)(3^{11}-1)/2.
\end{array}
\]
The first value is the number of totally isotropic planes in the
split 24-dimensional quadratic space over $\mathbb F_3$. These
equalities corroborate the classification and the shell counts; they do
not replace orbit-completeness or stabilizer proofs.

\section{The defining intersection and stability of the residue action}
We now make explicit which lattices carry the action in
Theorem~\ref{thm:residue-eta}. This stability condition is necessary:
equal rational representations do not themselves imply that a given
integral lattice is preserved.

Use the Euclidean model $A_5^4\perp D_4$ in four sum-zero six-spaces
and one four-space, with
$\lambda=(5,-1,-1,-1,-1,-1)/6$ and
$\omega_j=\sum_{i=1}^j e_i-(j/6)\sum_{i=1}^6e_i$.
The Niemeier lattice $N$ is obtained by adjoining the glue vectors
\[
 (3,3,3,3;0),\ (3,3,0,0;(1,0,0,0)),\
 (3,0,3,0;(\tfrac12,\tfrac12,\tfrac12,\tfrac12)),\
 (2,2,2,0;0),\ (2,4,0,2;0),
\]
where the first four entries multiply $\lambda$. Put
\[
 \rho=((5,3,1,-1,-3,-5)/2)^4\oplus(3,2,1,0),\qquad
 \alpha=(e_1-e_2,0,0,0;0),
\]
and, for $d\mid6$, define
\[
 K_d=\{x\in N:(\rho,x)\in d\mathbb Z\},\qquad
 N_d=K_d+\mathbb Z\bigl(\rho/d-(6/d)\alpha\bigr).
\]
Let $\Lambda_0=N_6$ and $J=\Lambda_0\cap c^2\Lambda_0$.

\begin{theorem}\label{thm:residue-stability}
The full six-cycle $c$ preserves $J$. Among its eight even unimodular
extensions, two have root system $A_2^{12}$ and are individually
preserved by $c$. The other six are Leech lattices, and $c$ permutes
them in two cycles of length three. Consequently:
\begin{enumerate}
\item The full group $S=\langle289\rangle$ acts on $J$ and on both
$A_2^{12}$ extensions; its Frame shapes and eta products are exactly
those in Theorem~\ref{thm:residue-eta}.
\item The subgroup preserving any one of the six Leech extensions
is $\langle289^3\rangle=\{1,169\}$. On such a fixed extension its
two eta products are $\Delta$ and $\eta(2\tau)^{12}$.
\end{enumerate}
For the defining intersection,
\begin{equation}\label{eq:theta-J-exact}
 \Theta_J=\frac{7381}{66430}E(\tau)
 +\frac{59049}{66430}E(3\tau)
 -\frac{264069}{11747}\Delta(\tau)
 -\frac{103276701}{11747}\Delta(3\tau)
 +\frac{14823}{1241}F(\tau).
\end{equation}
Every one of its eight unimodular extensions has $z=0$ in the same
five-dimensional basis. Thus the $F$ coordinate of
$\Theta_M-\Theta_J$ is $-14823/1241$ for every such extension $M$.
\end{theorem}

\begin{proof}
All lattice spans can be computed as integer Hermite forms after
multiplication of ambient coordinates by twelve. The five glue
classes generate an isotropic subgroup of order 72; hence the
determinant of $N$ is $6^4\cdot4/72^2=1$. Its evenness follows
from the integral pairings and half norms of the displayed vectors.
Since $(\rho,\alpha)=1$ and $\rho^2=84$, the map defining $K_d$
is onto. Its index is $d$, and the adjoining vector has norm
$4(6/d)^2$, pairs integrally with $K_d$, and has exact order $d$
modulo $K_d$. Thus $N_d$ is even unimodular.

The implementation forms intersections by the identity
$A\cap B=(A^\vee+B^\vee)^\vee$, using exact inverses and Hermite
forms. It verifies determinant 81 for $J$, equality of the three
pairwise intersections of $\Lambda_0,c^2\Lambda_0,c^4\Lambda_0$,
and invariance under $c^2$. There is also a short direct check for
$c^3$. Let $\beta$ be four copies of $\omega_3$, with zero $D_4$
component. Then $\beta\in N$ by its glue class $(3,3,3,3;0)$, and
$c^3\rho-\rho=-6\beta$. It follows that $c^3K_6=K_6$.
Writing $v=\rho/6-\alpha$, we have
$c^3v-v=-\beta-c^3\alpha+\alpha\in K_6$:
its pairing with $\rho$ is $-18-1+1=-18$.
Thus $c^3\Lambda_0=\Lambda_0$, and, as $c$ commutes with $c^2$,
also $c^3J=J$. Together with $c^2J=J$ this gives $cJ=J$.
The certificate additionally supplies the integral matrix of $c$
on a basis of $J$ and checks the full isometry identity.

Set $u=(\omega_2,\omega_4,0,\omega_2;0)$ and
$(e_1,e_2,e_3,e_4)=(u,v,c^2u,c^2v)$. Exact reduction shows that
their classes form a basis of $J^\vee/J\simeq\mathbb F_3^4$,
with Gram matrix and quadratic form
\[
 \begin{pmatrix}
 4&2&-2&-4/3\\2&4&-2/3&-2/3\\
 -2&-2/3&4&2\\-4/3&-2/3&2&4
 \end{pmatrix},\qquad
 Q(a,b,c,d)=\frac{2ad+bc+bd}{3}\pmod{\mathbb Z}.
\]
The matrix of $c$ on row coordinates in this quotient is
\[
 A=\begin{pmatrix}
 2&0&2&0\\1&2&2&2\\1&0&0&0\\0&1&0&0
 \end{pmatrix}\pmod3.
\]
In particular $A^3=1$. Even unimodular extensions are in bijection
with maximal totally isotropic planes of this split four-space.
Enumeration of its 81 vectors gives eight such planes. In the
certificate ordering, their root counts and $c$ permutation are
\[
 (0,0,72,0,0,72,0,0),\qquad
 (0\ 3\ 4)(1\ 7\ 6),
\]
with 2 and 5 fixed. The root counts come from exhaustive norm-two
enumeration. The rank-24 even unimodular classification identifies
the six rootless lattices as Leech lattices and the two with 72 roots
as $A_2^{12}$. This proves all the stability assertions.

Finally, independent enumeration for $J$ gives
\[
 \Theta_J=1+23058q^2+1849032q^3+44291124q^4+O(q^5),
 \qquad r_{3J^\vee}(6)=0.
\]
The norm-eight enumeration visited 261581721 recursion nodes and
used no floating-point inclusion decisions. Substitution into
Theorem~\ref{thm:universal-theta-coordinates} gives
\eqref{eq:theta-J-exact}; the additional $q^3,q^4$ counts also
verify it independently by the Sturm bound four.
For any even unimodular rank-24 lattice $M$, its level-one theta
series is $E+(r_M(2)-C)\Delta$. The six Leech cases therefore have
coordinates $(1,0,-C,0,0)$, and the two rooted cases have
$(1,0,72-C,0,0)$. The statements about $z$ follow.
\end{proof}

\section{Shell Dirichlet series and what Dirichlet GRH controls}
For $\Theta_J=1+\sum_{n\ge1}r_J(2n)q^n$, put
$D_J(s)=\sum_{n\ge1}r_J(2n)n^{-s}$. All formulas in this section
first hold for $\Re s>12$ and then continue meromorphically.
The decomposition above gives
\begin{equation}\label{eq:J-shell-series}
 D_J(s)=C(a+b3^{-s})\zeta(s)\zeta(s-11)
       +(x+y3^{-s})L(\Delta,s)+zL(F,s).
\end{equation}
Let $\chi$ be the nontrivial character modulo three and let $\chi_0$
be the principal character modulo three, both extended by zero on
multiples of three. Coefficientwise twisting gives
\[
 D_{J,\chi}(s)=aC L(s,\chi)L(s-11,\chi)
                  +xL(\Delta\otimes\chi,s)
                  +zL(F\otimes\chi,s),
\]
where the two terms supported on multiples of three disappear.
For the principal twist,
\begin{align*}
 D_{J,\chi_0}(s)={}&aC(1-3^{-s})(1-3^{11-s})
                        \zeta(s)\zeta(s-11)\\
 &+x(1-252\,3^{-s}+3^{11-2s})L(\Delta,s)
   +z(1+243\,3^{-s})L(F,s).
\end{align*}
These finite factors use $\tau(3)=252$ and $a_3(F)=-243$.

Dirichlet GRH locates the nontrivial zeros of the Dirichlet
$L$ factors, so the nontrivial zeros contributed by those factors
lie on $\Re s=1/2$ or $23/2$. It says nothing by itself about
the cuspidal factors or the zeros of the sums above. There are also
elementary zeros that must not be discarded. In particular the
untwisted Eisenstein summand has zeros at
\[
 s=\log_3\!\left(\frac{59049}{7381}\right)
       +\frac{(2j+1)\pi i}{\log3},\qquad j\in\mathbb Z,
\]
coming from $a+b3^{-s}$. The principal twist has additional finite-factor
zeros on $\Re s=0$ and $\Re s=11$. These assertions concern the
specified summands: adding cusp-form $L$ series can move or remove
such zeros. No transfer of a Riemann-hypothesis statement to
$D_J$ or its twists follows from linear decomposition alone.
